\documentclass[11pt]{article}

\usepackage[margin=1in]{geometry}
\usepackage{amsmath}
\usepackage{amssymb}
\usepackage{booktabs}
\usepackage{graphicx}
\usepackage{float}
\usepackage{array}
\usepackage{longtable}
\usepackage{hyperref}

\title{\textbf{Final Model Report: NYC Uber/Lyft Hourly Demand Forecasting}}
\author{Tommy Yao}
\date{August-September 2026}

\begin{document}

\maketitle

% ============================================================
\begin{abstract}

This report presents the final demand forecasting model for hourly
Uber and Lyft demand across New York City pickup zones. The final
framework combines a primary XGBoost demand forecasting model with
a residual correction model designed specifically for high-demand
observations.

The final model achieves an overall mean absolute error (MAE) of
8.4533, root mean squared error (RMSE) of 15.9376, and coefficient
of determination ($R^2$) of 0.9516. Model performance remains
particularly strong for observations between 50 and 500 trips per
zone-company-hour. However, extreme demand above 500 trips remains
the primary source of forecasting error, with a normalized MAE of
12\% and an underprediction rate of 53\%.

The results suggest that the model captures the majority of regular
temporal demand patterns effectively, while extreme demand events
remain more difficult to predict because of their relatively low
frequency and larger deviations from historical demand patterns.

\end{abstract}




% ============================================================

\section{Correction Model Training Data}

The high-demand validation sample contains 2,879 observations.

\begin{table}[H]
\centering
\caption{High-Demand Correction Data}
\begin{tabular}{lr}
\toprule
Statistic & Value \\
\midrule
High-demand observations & 2,879 \\
Average actual demand & 584.7100 \\
Average main-model prediction & 538.8819 \\
Average correction target & 45.8281 \\
\bottomrule
\end{tabular}
\end{table}

The positive average correction target indicates that the main model
systematically underestimated demand within the high-demand correction
sample. On average, actual demand exceeded the main-model prediction
by approximately 45.83 trips.


% ============================================================
\section{Overall Final Model Performance}

The final model produces the following out-of-sample performance:

\begin{table}[H]
\centering
\caption{Final Model Performance}
\begin{tabular}{lr}
\toprule
Metric & Value \\
\midrule
MAE & 8.4563 \\
RMSE & 15.7416 \\
$R^2$ & 0.9516 \\
\bottomrule
\end{tabular}
\end{table}

The MAE indicates that the model's prediction differs from observed
hourly demand by approximately 8.45 trips on average.

The RMSE is larger than the MAE because large forecasting errors
receive greater penalties under the squared-error formulation.
This difference indicates that a relatively small number of
high-error observations remain in the validation data.

The $R^2$ of 0.9504 indicates that approximately 95.04\% of the
variation in demand within the evaluation sample is explained by
the model.


% ============================================================
\section{Forecasting Error by Demand Level}

Model performance varies considerably across demand regimes.

Define normalized MAE as

\begin{equation}
\mathrm{NMAE}
=
\frac{
\sum_{i=1}^{n}
\left|
D_i-\widehat{D}_i
\right|
}{
\sum_{i=1}^{n} D_i
}.
\end{equation}

\begin{table}[htbp]
\centering
\caption{Prediction Direction by Demand Level}
\label{tab:prediction_direction}
\begin{tabular}{lcc}
\hline
\textbf{Demand Level} & \textbf{Overprediction} & \textbf{Underprediction} \\
\hline
0--5     & 0.729496 & 0.270504 \\
5--10    & 0.583647 & 0.416353 \\
10--50   & 0.551257 & 0.448743 \\
50--100  & 0.496336 & 0.503664 \\
100--200 & 0.407970 & 0.592030 \\
200--500 & 0.330779 & 0.669221 \\
500+     & 0.374809 & 0.625191 \\
\hline
\end{tabular}
\end{table}


\subsection{Interpretation}

The model performs particularly well in the medium- and high-demand
ranges between 50 and 500 trips. Normalized MAE is approximately
11--13\% across these observations.

The very high normalized error for the 0--5 demand category should
be interpreted carefully. Because the denominator is extremely
small, even an absolute error of only one or two trips can produce
a large normalized percentage error. The absolute MAE in this
category remains only 1.95 trips. In additional to this, small
demand does not have a large business impact because the demand
level is relatively small and it is less lucrative.

Extreme demand above 500 trips presents an improvement when I used
the additional correction modeol using XGboost.Average actual demand is 
approximately 634 trips, while the average
prediction is approximately 595 trips. Therefore, the remaining
error is somehow acceptable but there can be a room for advanced
improvenment.


% ============================================================
\section{Overprediction and Underprediction}

The direction of forecasting error provides additional information
beyond absolute-error metrics.

\begin{table}[H]
\centering
\caption{Prediction Direction by Demand Level}
\begin{tabular}{lrr}
\toprule
Demand Level &
Overprediction &
Underprediction \\
\midrule
0--5    & 72.95\% & 27.05\% \\
5--10   & 58.36\% & 41.64\% \\
10--50  & 55.13\% & 44.87\% \\
50--100 & 49.63\% & 50.37\% \\
100--200 & 40.80\% & 59.20\% \\
200--500 & 33.34\% & 66.92\% \\
500+     & 37.48\% & 62.51\% \\
\bottomrule
\end{tabular}
\end{table}

A clear directional pattern is observed. The model tends to
overpredict very low demand and increasingly underpredict demand as
the actual demand level rises.

After the correction model is applied,For demand above 500 trips, 
approximately 53.39\% of observationsare underpredicted. This 
indicates that extreme demand remains the
primary regime in which systematic prediction bias persists.


% ============================================================
\section{Extreme Demand Analysis}

For observations with actual demand exceeding 500 trips:

\begin{table}[H]
\centering
\caption{Performance for Demand Above 450 Trips}
\begin{tabular}{lr}
\toprule
Statistic & Value \\
\midrule
Observations & 2879 \\
Average Actual Demand & 584.7 \\
Average Prediction & 538.88 \\
MAE & 75.01 \\
NMAE & 12.8274\% \\
Underprediction Rate & 53.58\% \\
\bottomrule
\end{tabular}
\end{table}

Despite the residual correction mechanism, extreme demand remains
more difficult to forecast than regular demand. Nevertheless, an
NMAE of 16.49\% indicates that the predictions retain substantial
information about the magnitude of these high-demand periods.


% ============================================================
\section{Spatial Analysis of Extreme Demand}

Extreme-demand observations are concentrated in a relatively small
number of major transportation and commercial zones.

\begin{table}[htbp]
\centering
\caption{High-Demand Prediction Performance by Pickup Zone}
\label{tab:high_demand_zone}
\begin{tabular}{lrrrrr}
  \hline
\textbf{Pickup Zone} &
\textbf{Observations} &
\textbf{Avg. Demand} &
\textbf{Avg. Prediction} &
\textbf{Avg. Residual} &
\textbf{MAE} \\
\hline
LaGuardia Airport          & 597 & 667.74 & 629.06 & 38.68  & 92.18 \\
JFK Airport                & 195 & 602.29 & 586.19 & 16.10  & 58.17 \\
Midtown Center             & 139 & 701.17 & 692.96 & 8.20   & 49.80 \\
East Village               & 134 & 719.54 & 682.65 & 36.88  & 82.65 \\
Lower East Side            & 101 & 684.72 & 647.27 & 37.45  & 86.55 \\
TriBeCa/Civic Center       & 82  & 567.04 & 507.94 & 59.09  & 63.99 \\
Union Sq                   & 72  & 554.67 & 529.87 & 24.79  & 41.92 \\
West Chelsea/Hudson Yards  & 70  & 587.94 & 538.47 & 49.47  & 59.00 \\
Williamsburg (North Side)  & 69  & 595.99 & 561.85 & 34.13  & 60.62 \\
Crown Heights North        & 69  & 578.29 & 528.25 & 50.04  & 62.32 \\
Times Sq/Theatre District  & 64  & 584.36 & 586.35 & -1.99  & 46.60 \\
Bushwick South             & 62  & 654.94 & 603.80 & 51.13  & 83.33 \\
West Village               & 53  & 593.49 & 566.07 & 27.42  & 47.54 \\
Greenpoint                 & 47  & 571.04 & 541.33 & 29.71  & 63.17 \\
Midtown East               & 41  & 609.83 & 579.58 & 30.25  & 46.98 \\
Bushwick North             & 35  & 621.49 & 548.25 & 73.23  & 97.36 \\
\hline
\end{tabular}
\end{table}

LaGuardia Airport contains the largest number of extreme-demand
observations, followed by JFK Airport. Extreme demand is also
concentrated in major Manhattan commercial and nightlife areas,
including Midtown Center, East Village, Lower East Side, and
West Chelsea/Hudson Yards.

The spatial results also reveal substantial heterogeneity in
forecasting difficulty. For example, average residuals are
considerably larger in several nightlife and residential-commercial
zones than at JFK Airport. This suggests that extreme demand may be
generated by different mechanisms across zones.


% ============================================================
\section{Holiday Performance}

\begin{table}[H]
\centering
\caption{Holiday and Non-Holiday Forecast Performance}
\begin{tabular}{lrrrr}
\toprule
Holiday &
Observations &
Avg. Demand &
Avg. Prediction &
MAE \\
\midrule
No
& 684,774 & 58.912 & 58.629 & 8.384 \\

Yes
& 35,287 & 50.361 & 50.039 & 9.790 \\
\bottomrule
\end{tabular}
\end{table}

Average predictions remain close to average observed demand for both
holiday and non-holiday periods. However, holiday observations have
a higher MAE of 9.79 compared with 8.38 for non-holiday observations.

This indicates that holiday demand contains greater observation-level
variation even though the model remains approximately unbiased at the
aggregate level.


% ============================================================
\section{Feature Importance}

The XGBoost feature importance results are shown below.

\begin{table}[H]
\centering
\caption{Final Model Feature Importance}
\begin{tabular}{lr}
\toprule
Feature & Importance \\
\midrule
Lag 168 Hours & 0.4357 \\
Lag 1 Hour & 0.3974 \\
Lag 24 Hours & 0.0770 \\
Rolling Mean 3 Hours & 0.0364 \\
Zone-Hour Interaction & 0.0147 \\
Hour & 0.0109 \\
Day of Week & 0.0036 \\
Pickup Zone & 0.0032 \\
Rolling Maximum 24 Hours & 0.0032 \\
Holiday Indicator & 0.0028 \\
Rolling Mean 24 Hours & 0.0025 \\
Weekend Indicator & 0.0020 \\
Rain Category & 0.0018 \\
Rolling Mean 168 Hours & 0.0016 \\
Temperature & 0.0016 \\
Snow Indicator & 0.0016 \\
Month & 0.0015 \\
Wind Category & 0.0014 \\
Company & 0.0012 \\
\bottomrule
\end{tabular}
\end{table}


\subsection{Interpretation of Feature Importance}

Historical demand dominates the predictive structure of the model.
The two most important predictors are demand from the same hour one
week earlier and demand from the immediately preceding hour.

Together,

\begin{equation}
0.4357 + 0.3974 \approx 0.8331,
\end{equation}

meaning that these two variables account for approximately 83.3\%
of the reported XGBoost feature importance.

The importance of $\mathrm{Lag168}$ demonstrates strong weekly
periodicity in ride demand, while $\mathrm{Lag1}$ captures
short-term demand persistence. The importance of $\mathrm{Lag24}$
provides additional evidence of daily seasonality.

Weather variables have relatively low individual feature importance.
This does not necessarily imply that weather has no effect on ride
demand. Instead, much of its predictive contribution may be smaller
than, conditional on, or partially captured by recent historical
demand patterns.


% ============================================================
\section{Business Interpretation}

The model produces three findings that are particularly relevant
for operational demand forecasting.

\begin{enumerate}

\item \textbf{Regular demand is highly predictable.}

For demand between 50 and 500 trips per zone-company-hour, normalized
MAE remains approximately 11--13\%. This suggests that historical
demand patterns provide a strong operational baseline for forecasting
normal and moderately high demand.

\item \textbf{Extreme demand remains systematically underestimated.}

For demand above 500 trips, the model underpredicts approximately
81.39\% of observations. Therefore, an operational system relying
only on point forecasts should account for additional uncertainty
during extreme-demand periods.

\item \textbf{Extreme-demand errors are spatially concentrated.}

Airports, major Manhattan commercial areas, and several nightlife
zones account for a substantial portion of observations above
500 trips. Consequently, further model improvements could focus on
zone-specific extreme-demand behavior rather than applying a uniform
correction across all New York City zones.

\end{enumerate}


% ============================================================
\section{Model Limitations}

Several limitations remain in the final forecasting framework.

First, extreme-demand observations are relatively rare compared with
normal-demand observations. Only 2848 validation observations exceed
500 trips, making this regime substantially more difficult for the
model to learn.

Second, the current feature set relies heavily on historical demand.
Unexpected events such as concerts, major sporting events, flight
disruptions, transit interruptions, or unusual local activity may
generate demand shocks that cannot be fully inferred from lagged
demand.

Third, standard XGBoost feature importance measures predictive usage
within the fitted model and should not be interpreted as causal
effects. In particular, the relatively low importance assigned to
weather variables does not demonstrate that weather has no causal
effect on ride demand.

Finally, the residual correction model reduces systematic high-demand
bias but does not eliminate it. The remaining 53.58\% underprediction
rate for demand above 500 indicates that extreme-demand forecasting
remains the primary opportunity for future improvement.


% ============================================================
\section{Potential Future Improvements}

Future extensions could incorporate additional predictors related to
extreme-demand events, including

\begin{itemize}

\item major concerts and sporting events,
\item airport flight arrival and departure volume,
\item public transportation disruptions,
\item major holidays and special events,
\item zone-specific high-demand correction models,
\item interaction effects between location, hour, and weather, and
\item probabilistic prediction intervals for operational uncertainty.

\end{itemize}

A particularly useful extension would be to model extreme demand
separately for airports, Manhattan commercial zones, and nightlife
areas. The spatial analysis suggests that the mechanisms generating
extreme demand differ substantially across these locations.


% ============================================================
\section{Final Conclusion}

The final model demonstrates strong overall forecasting performance, but systematic underprediction remains during extreme-demand periods.
 In particular, demand above 500 trips is underpredicted in approximately 53.59\% of observations. However, these extreme-demand observations
  are not distributed uniformly across New York City. They are concentrated in several high-volume locations, including LaGuardia Airport, 
  JFK Airport, Midtown, East Village, and other major commercial or nightlife zones.

Rather than uniformly increasing predictions across all zones, a more practical operational strategy is to introduce a \textbf{targeted demand buffer} for locations with historically persistent high-demand underprediction. 
For a selected set of high-risk zones $\mathcal{Z}_{H}$, an operational forecast can be expressed as

\begin{equation}
\widehat{D}_{z,c,t}^{\mathrm{operational}}
=
\widehat{D}_{z,c,t}^{\mathrm{final}}
+
I(z \in \mathcal{Z}_{H})B_{z,t},
\end{equation}

where $B_{z,t}$ represents an additional demand buffer determined from historical residual patterns for zone $z$ and potentially varying by time of day.

This approach provides additional capacity protection in locations where extreme demand is repeatedly underestimated without unnecessarily increasing predictions throughout the entire system.

At the same time, the buffer should remain conservative. Excessive upward adjustment may create \textbf{overprediction}, which can result in too many drivers being directed toward locations where realized demand is insufficient. This can increase driver idle time, empty vehicle mileage, congestion, and inefficient allocation of available supply. Therefore, completely eliminating underprediction should not necessarily be the operational objective.

A more appropriate objective is to balance the asymmetric costs of forecasting errors. Moderate underprediction can often be addressed through real-time driver repositioning and supply responses, while substantial overprediction may proactively allocate resources to demand that never materializes. Consequently, the recommended strategy is to preserve the current model as the primary forecast while applying \textbf{small, zone-specific buffers only to locations and periods with persistent extreme-demand underprediction}.

Overall, the model should therefore be viewed not only as a point forecasting system, but as a decision-support framework in which statistical predictions are combined with targeted operational adjustments. This approach reduces exposure to extreme-demand shortages while limiting the operational cost associated with unnecessary overprediction.


\end{document}